\begin{abstract}

The Prescription Drug User Fee Act (PDUFA) of 1992 introduced industry user fees and performance targets that substantially accelerated the FDA drug approval process. While existing research primarily evaluates the implications of faster regulatory review for post-approval drug safety, less attention has been given to how regulatory acceleration may affect the composition of drugs entering the market. This paper examines whether the introduction of PDUFA altered the composition of FDA-approved drugs in the United States, focusing on pharmaceuticals classified as controlled substances under the Controlled Substances Act. Using data on FDA drug approvals, regulatory review pathways, and drug scheduling status, the analysis investigates whether the share of approved drugs with abuse potential changed following the implementation of user-fee–funded accelerated review. The study also examines whether drugs classified as controlled substances are more likely to receive expedited regulatory designations such as priority review or other accelerated programs. By studying how regulatory reforms influence the composition of approved pharmaceuticals, this project contributes to the literature on pharmaceutical regulation, innovation incentives, and the unintended consequences of health policy.

\end{abstract}
